\documentclass[11pt,a4paper]{article}
\usepackage[T1]{fontenc}
\usepackage{lmodern}
\usepackage[margin=25mm]{geometry}
\usepackage{amsmath,amssymb,amsthm,mathtools}
\usepackage{enumitem,microtype,xcolor,booktabs}
\usepackage[hyperpageref]{backref}
\usepackage[pagebackref]{hyperref}
\hypersetup{
    colorlinks=true,
    linkcolor=red,
    filecolor=blue,
    urlcolor=orange,
    citecolor=violet
}
\renewcommand*{\backref}[1]{}
\renewcommand*{\backrefalt}[4]{
  \ifcase #1 Not cited.
  \or Cited on page #2.
  \else Cited on pages #2.
  \fi}
\definecolor{darkblue}{rgb}{0,0,0.7}
\newcommand{\defn}[1]{\textsl{\color{darkblue} #1}}
\setlist[enumerate]{leftmargin=*,itemsep=.5em,topsep=.4em}
\setlength{\emergencystretch}{2em}
\newcommand{\D}{\mathcal D}
\renewcommand{\H}{\mathcal H}
\newcommand{\ZZ}{\mathbb{Z}}
\newcommand{\op}{\mathrm{op}}
\newcommand{\RR}{\mathbf{R}}
\newcommand{\std}{\Delta}
\newcommand{\costd}{\nabla}
\newcommand{\dual}{\mathbb D}
\newcommand{\dgHom}{\mathcal{H}\!om}
\DeclareMathOperator{\fd}{fd}
\DeclareMathOperator{\Hom}{Hom}
\DeclareMathOperator{\RHom}{\RR Hom}
\DeclareMathOperator{\REnd}{\RR End}
\DeclareMathOperator{\add}{add}
\DeclareMathOperator{\thick}{thick}
\DeclareMathOperator{\per}{per}
\DeclareMathOperator{\Filt}{Filt}
\let\mod\relax
\DeclareMathOperator{\mod}{mod}
\theoremstyle{plain}
\newtheorem{theorem}{Theorem}[section]
\newtheorem{proposition}[theorem]{Proposition}
\newtheorem{corollary}[theorem]{Corollary}
\theoremstyle{definition}
\newtheorem{definition}[theorem]{Definition}
\newtheorem{example}[theorem]{Example}
\numberwithin{equation}{section}
\title{$m$-reflexive qh dgas}
\author{}
\date{}
\begin{document}
\maketitle
\begin{abstract}
We study exceptional standard objects in the $m$-extended module category
whose left-dual objects have cohomology in degrees $0,\ldots,m-1$.
The definition specialises to classical quasi-heredity for $m=1$.
For every ordering of the vertices of a finite acyclic graded quiver,
we construct suitable standard dg modules and prove fullness in both
$\per A$ and $\D_{\fd}(A)$. The construction allows a non-zero differential;
its cohomological bounds are controlled by the degrees of paths.
\end{abstract}

\section{Definition and classical specialisation}\label{sec:definition}

Fix an algebraically closed field $\Bbbk$ and an integer $m\geq1$.
All dg modules are right modules, and all complexes are cochain complexes:
$d_X:X^q\to X^{q+1}$ and $H^q(X[s])=H^{q+s}(X)$.
Let $A$ be a proper connective dg algebra satisfying
\begin{equation}\label{eq:standing}
 \dim_{\Bbbk}H^*(A)<\infty,\qquad
 H^q(A)=0\quad(q\notin[1-m,0]).
\end{equation}
Assume that $H^0(A)$ is basic and that closed degree-zero orthogonal
idempotents $e_1,\ldots,e_n$ sum to $1$, with classes forming a complete
set of primitive idempotents of $H^0(A)$. Set
\begin{equation}\label{eq:categories}
 \begin{gathered}
 P_v=e_vA,\qquad
 \D_{\fd}(A)=\{X\in\D(A):\dim_{\Bbbk}H^*(X)<\infty\},\\
 \per A=\thick_{\D(A)}(A),\qquad
 \H_A^m=\D_{\fd}(A)^{[1-m,0]}.
 \end{gathered}
\end{equation}
For a class $\mathcal S\subseteq\D_{\fd}(A)$,
$\thick_{\D_{\fd}(A)}(\mathcal S)$ denotes closure under shifts, cones,
and direct summands. Properness gives $\per A\subseteq\D_{\fd}(A)$.

The category $\H_A^m$ is extriangulated, with
\begin{equation}\label{eq:extension}
 \mathbb E_{\H_A^m}(X,Y)=\Hom_{\D_{\fd}(A)}(X,Y[1]),
 \qquad \mathcal P(\H_A^m)=\add A
 \quad(X,Y\in\H_A^m)
\end{equation}
\cite[Corollary 2.3]{MP25}. A conflation
$K\to P\to X\dashrightarrow$ in $\H_A^m$ is induced by a triangle
$K\to P\to X\to K[1]$ in $\D_{\fd}(A)$ with $K,P,X\in\H_A^m$.
We write $\Filt_{\H_A^m}(\mathcal S)$ for the finite extension closure
of $\mathcal S$ in $\H_A^m$, including zero and finite direct sums.

An object $E\in\D_{\fd}(A)$ is \defn{exceptional} if
\begin{equation}\label{eq:exceptional}
 \Hom_{\D_{\fd}(A)}(E,E[t])\cong
 \begin{cases}
 \Bbbk,&t=0,\\
 0,&t\neq0
 \end{cases}
\qquad(t\in\ZZ).
\end{equation}
An ordered tuple $(E_1,\ldots,E_n)$ of exceptional objects is an
\defn{exceptional sequence} if
\begin{equation}\label{eq:exceptional-sequence}
 \Hom_{\D_{\fd}(A)}(E_i,E_j[t])=0
 \qquad(1\leq j<i\leq n,\ t\in\ZZ).
\end{equation}
The indices $i,j$ in \eqref{eq:exceptional-sequence} denote positions
in the sequence.
For exceptional objects $\std_1,\ldots,\std_n$, a permutation
$\sigma\in S_n$ lists the vertex labels in the standard sequence
\begin{equation}\label{eq:standard-sequence}
 \boldsymbol\std^\sigma
 =(\std_{\sigma(1)},\ldots,\std_{\sigma(n)}).
\end{equation}
Thus $\sigma^{-1}(v)$ is the position of $\std_v$ in
$\boldsymbol\std^\sigma$. Define the total order $\leq_\sigma$ on
the vertex labels $\{1,\ldots,n\}$ by
\begin{equation}\label{eq:order}
 u\leq_\sigma v\quad\Longleftrightarrow\quad
 \sigma^{-1}(u)\leq\sigma^{-1}(v).
\end{equation}
We write $u<_\sigma v$ if $u\leq_\sigma v$ and $u\neq v$.
The least and greatest weights are $\sigma(1)$ and $\sigma(n)$,
respectively. Thus smaller weights occur earlier in
\eqref{eq:standard-sequence}. By \eqref{eq:exceptional-sequence},
the sequence $\boldsymbol\std^\sigma$ is exceptional precisely when
\begin{equation}\label{eq:semiorthogonal}
 \Hom_{\D_{\fd}(A)}(\std_w,\std_v[t])=0
 \qquad(v<_\sigma w,\ t\in\ZZ).
\end{equation}
The \defn{left-dual sequence} of $\boldsymbol\std^\sigma$ is an
exceptional sequence $(\costd_{\sigma(n)},\ldots,\costd_{\sigma(1)})$,
listed in decreasing order for $\leq_\sigma$, whose terms satisfy
\begin{equation}\label{eq:left-dual-membership}
 \costd_v\in\thick_{\D_{\fd}(A)}(\std_1,\ldots,\std_n)
 \qquad(1\leq v\leq n)
\end{equation}
and whose pairing with the standard objects is
\begin{equation}\label{eq:pairing}
 \Hom_{\D_{\fd}(A)}(\std_u,\costd_v[t])\cong
 \begin{cases}
 \Bbbk,&u=v,\ t=0,\\
 0,&\text{otherwise}
 \end{cases}
 \qquad(1\leq u,v\leq n,\ t\in\ZZ).
\end{equation}
We use the convention of \cite[\S2.3]{BB26}. The objects $\costd_v$,
when they exist, are determined up to isomorphism in
$\thick_{\D_{\fd}(A)}(\std_1,\ldots,\std_n)$ by \eqref{eq:pairing}.

\begin{definition}[$m$-reflexive qh dga]\label{def:reflexive}
An algebra $A$ satisfying \eqref{eq:standing} is an
\defn{$m$-reflexive qh dga with respect to $\sigma$} if there exist
objects $\std_v\in\H_A^m$ such that:
\begin{enumerate}[label=\textup{(R\arabic*)},ref=R\arabic*]
 \item\label{ax:exceptional}
 $\boldsymbol\std^\sigma$ satisfies
 \eqref{eq:exceptional} and \eqref{eq:semiorthogonal};
 \item\label{ax:presentation}
 for each vertex $v$, there is a conflation
 \begin{equation}\label{eq:presentation}
 K_v\longrightarrow P_v\xrightarrow{p_v}\std_v\dashrightarrow
 \qquad\text{in }\H_A^m;
 \end{equation}
 \item\label{ax:costandard}
 the left-dual sequence of $\boldsymbol\std^\sigma$ exists, with
 \begin{equation}\label{eq:costandard-window}
 \costd_v\in\D_{\fd}(A)^{[0,m-1]}
 \qquad(1\leq v\leq n);
 \end{equation}
 \item\label{ax:full}
 $\boldsymbol\std^\sigma$ is full:
 \begin{equation}\label{eq:full}
 \thick_{\D_{\fd}(A)}(\std_1,\ldots,\std_n)=\D_{\fd}(A).
 \end{equation}
\end{enumerate}
\end{definition}

In Definition~\ref{def:reflexive}, $\boldsymbol\std^\sigma$ is the chosen
exceptional sequence; condition~\ref{ax:costandard} concerns the
left-dual objects $\costd_v$ specified by \eqref{eq:pairing}.
Equation~\eqref{eq:pairing} implies $H^0(\std_v)\neq0$:
otherwise $\std_v\in\D_{\fd}(A)^{\leq-1}$ and
$\costd_v\in\D_{\fd}(A)^{\geq0}$ would give
$\Hom_{\D_{\fd}(A)}(\std_v,\costd_v)=0$.
The epimorphism
$H^0(p_v):H^0(P_v)\to H^0(\std_v)$ is therefore a projective cover.
In particular, the top of $H^0(\std_v)$ is the simple $H^0(A)$-module
at vertex $v$.

\begin{proposition}[Classical specialisation]\label{prop:classical}
Suppose $H^q(A)=0$ for $q\neq0$, and put $B=H^0(A)$.
The following conditions are equivalent:
\begin{enumerate}[label=\textup{(\roman*)}]
 \item $A$ is a $1$-reflexive qh dga with respect to $\sigma$;
 \item $B$ is quasi-hereditary with respect to $\sigma$.
\end{enumerate}
\end{proposition}
\begin{proof}
The quasi-isomorphism $A\simeq B$ identifies $\H_A^1$ with $\mod B$.

\noindent$\underline{(\mathrm{ii})\Rightarrow(\mathrm{i})}$.
The classical standard modules have projective-cover presentations
$P_v\twoheadrightarrow\std_v$ whose kernels are filtered by
$\std_w$ with $v<_\sigma w$ \cite{DR92}.
The classical standard and costandard sequences are full exceptional
sequences satisfying \eqref{eq:pairing}
\cite[Lemma 2.2]{BB26}. The classical projective-cover presentations
and standard and costandard sequences satisfy
conditions~\ref{ax:exceptional}--\ref{ax:full}.

\noindent$\underline{(\mathrm{i})\Rightarrow(\mathrm{ii})}$.
Conditions~\ref{ax:exceptional}, \ref{ax:costandard}, and \ref{ax:full}
give full exceptional sequences
$(\std_{\sigma(1)},\ldots,\std_{\sigma(n)})$ and
$(\costd_{\sigma(n)},\ldots,\costd_{\sigma(1)})$ in
$\D^b(\mod B)$, with all terms in $\mod B$.
By \cite[Theorem 3.8]{BB26}, the sequences
$\boldsymbol\std^\sigma$ and
$(\costd_{\sigma(n)},\ldots,\costd_{\sigma(1)})$
determine a highest weight structure on $\mod B$.
The projective covers $H^0(p_v)$ identify the highest weight structure's
standard-object labels with the vertices ordered by $\sigma$.
\end{proof}

\begin{example}[The order for two classical standards]\label{ex:classical-order}
Take $A=B=\Bbbk(1\xrightarrow{a}2)$ concentrated in degree zero,
with $e_1a=a=ae_2$, $m=1$, and $\sigma=\mathrm{id}$. Thus $1<_\sigma2$.
Let $S_v$ be the simple top of $P_v=e_vB$, for $v=1,2$. Then
\begin{equation}\label{eq:classical-order-objects}
 \std_1=S_1,\qquad \std_2=S_2=P_2,\qquad
 \costd_1=S_1,\qquad \costd_2=P_1.
\end{equation}
The projective presentation
\begin{equation}\label{eq:classical-order-extension}
 0\longrightarrow S_2\longrightarrow P_1\longrightarrow S_1
 \longrightarrow0
\end{equation}
gives
\begin{equation}\label{eq:classical-order-hom}
 \Hom_{\D_{\fd}(B)}(S_2,S_1[t])=0\quad(t\in\ZZ),\qquad
 \Hom_{\D_{\fd}(B)}(S_1,S_2[1])\cong\Bbbk.
\end{equation}
Consequently $(\std_1,\std_2)=(S_1,S_2)$ is full exceptional,
whereas $(\std_2,\std_1)$ is not exceptional. The left-dual sequence
of $(\std_1,\std_2)$ specified by \eqref{eq:pairing} is
$(\costd_2,\costd_1)=(P_1,S_1)$.
\end{example}

\section{Graded path dg algebras}\label{sec:paths}

Let $Q$ be a finite acyclic quiver with vertex set
$Q_0=\{1,\ldots,n\}$, where $n\geq1$.
Give each arrow an integer degree at most $0$.
For a path $p$, let $s(p)$, $t(p)$, and $|p|$ denote its source,
target, and degree. Trivial paths have degree $0$.
We concatenate paths from left to right, so
$e_ua=a=ae_v$ for an arrow $a:u\to v$.
Let
\begin{equation}\label{eq:path-dga}
 A=(\Bbbk Q,d_A),\qquad
 d_A:A^q\longrightarrow A^{q+1},\qquad d_A^2=0,
\end{equation}
where $d_A$ is a graded derivation. Since $A^1=0$, all $e_v$ are closed.
Define
\begin{equation}\label{eq:path-width}
 M=1-\min\{|p|:p\text{ is a path in }Q\}.
\end{equation}
Thus $A^q=0$ for $q\notin[1-M,0]$.

For complexes of $\Bbbk$-vector spaces, write
$\dual V=\Hom^\bullet_{\Bbbk}(V,\Bbbk)$, with
\[
 (\dual V)^q=\Hom_{\Bbbk}(V^{-q},\Bbbk),
 \qquad d_{\dual V}(\varphi)=(-1)^{q+1}\varphi d_V
 \quad(\varphi\in(\dual V)^q).
\]
For a left dg $A$-module $V$, the right action on $\dual V$ is
$(\varphi a)(x)=\varphi(ax)$.
For right dg modules $X,Y$, let $\dgHom_A(X,Y)$ be the cochain complex
of homogeneous $A$-linear maps, with
\[
 d(f)=d_Yf-(-1)^q f d_X
 \quad\bigl(f\in\dgHom_A(X,Y)^q\bigr).
\]
We write $\RHom_A(X,Y)$ for the corresponding derived Hom complex.

\begin{theorem}[Standards for every ordering]\label{thm:paths}
Let $A$ be the path dga \eqref{eq:path-dga}, and fix $\sigma\in S_n$.
For each vertex $v$, set
\begin{equation}\label{eq:path-objects}
 \begin{gathered}
 e_{>v}^{\sigma}=\sum_{v<_\sigma w}e_w,\qquad
 B_v=A/Ae_{>v}^{\sigma}A,\qquad
 K_v=e_vAe_{>v}^{\sigma}A,\\
 \std_v=e_vB_v=P_v/K_v,\qquad
 C_v=B_ve_v,\qquad
 \costd_v=\dual C_v.
 \end{gathered}
\end{equation}
Then $A$ is an $M$-reflexive qh dga with respect to $\sigma$, with
standard objects $\std_v$, left-dual objects $\costd_v$, and
conflations induced by
\begin{equation}\label{eq:path-short-exact}
 0\longrightarrow K_v\longrightarrow P_v
       \xrightarrow{p_v}\std_v\longrightarrow0.
\end{equation}
Moreover,
\begin{equation}\label{eq:path-full}
 \thick_{\D_{\fd}(A)}(\std_1,\ldots,\std_n)
       =\per A=\D_{\fd}(A).
\end{equation}
The objects in \eqref{eq:path-objects} also satisfy
Definition~\ref{def:reflexive} for every $m\geq M$.
\end{theorem}

\begin{proof}
The ideals $Ae_{>v}^{\sigma}A$ are dg ideals, because the idempotents
$e_w$ are closed. Thus every object in \eqref{eq:path-objects} is a dg
module, and \eqref{eq:path-short-exact} is a short exact sequence of dg
modules.

\begin{enumerate}[label=\textbf{\arabic*.},ref=\arabic*]
\item\label{step:kernels}
\textbf{Finite projective filtrations of the kernels.}
Let $\mathcal F_v$ be the set of paths
$p=(v=v_0\to v_1\to\cdots\to v_\ell)$ such that
\[
 v<_\sigma v_\ell,\qquad
 v_j\leq_\sigma v\quad(0\leq j<\ell).
\]
Every basis path of $K_v$ has a unique prefix in $\mathcal F_v$.
As graded right modules,
\begin{equation}\label{eq:graded-kernel}
 K_v=\bigoplus_{p\in\mathcal F_v}pA
       \cong\bigoplus_{p\in\mathcal F_v}P_{t(p)}[-|p|].
\end{equation}
The summands $pA$ need not be dg submodules.
Write the unique expansion
\begin{equation}\label{eq:generator-differential}
 \begin{gathered}
 d_A(p)=\sum_{p'\in\mathcal F_v}p'a_{p'p},\\
 a_{p'p}\in e_{t(p')}A^{|p|+1-|p'|}e_{t(p)}.
 \end{gathered}
\end{equation}
Every non-zero homogeneous coefficient $a_{p'p}$ satisfies
$|p'|=|p|+1-|a_{p'p}|\geq |p|+1$.
Consequently,
\begin{equation}\label{eq:kernel-filtration}
 F^hK_v=\bigoplus_{\substack{p\in\mathcal F_v\\|p|\geq h}}pA,
 \qquad
 F^hK_v/F^{h+1}K_v
 \cong\bigoplus_{\substack{p\in\mathcal F_v\\|p|=h}}P_{t(p)}[-h]
\end{equation}
is a finite filtration by dg submodules, with $h\in\ZZ$.
Indeed, modulo $F^{h+1}K_v$,
$d_A(px)=(-1)^h p\,d_A(x)$ for $|p|=h$, which is the differential
on $P_{t(p)}[-h]$. Each vertex $t(p)$ satisfies $v<_\sigma t(p)$.

Filtration \eqref{eq:kernel-filtration} makes $K_v$ perfect and
homotopically projective: $\dgHom_A(K_v,-)$ preserves
quasi-isomorphisms. Sequence \eqref{eq:path-short-exact} then makes
$\std_v$ perfect. Descending induction in $<_\sigma$, using
\eqref{eq:kernel-filtration} and \eqref{eq:path-short-exact}, gives
\begin{equation}\label{eq:perfect-generation}
 P_v\in\thick_{\D_{\fd}(A)}(\std_v,\std_w:v<_\sigma w),
 \qquad
 \thick_{\D_{\fd}(A)}(\std_1,\ldots,\std_n)=\per A.
\end{equation}

\item\label{step:exceptional}
\textbf{Exceptionality of the standard sequence.}
Suppose $v\leq_\sigma u$. Every projective factor of $K_u$ in
\eqref{eq:kernel-filtration} is a shift of $P_w$ with
$u<_\sigma w$. Since $\std_v e_w=0$,
\[
 \RHom_A(K_u,\std_v)=0,\qquad
 \RHom_A(P_u,\std_v)\simeq\std_v e_u
 =
 \begin{cases}\Bbbk,&u=v,\\0,&v<_\sigma u.\end{cases}
\]
The triangle induced by \eqref{eq:path-short-exact} yields
\eqref{eq:exceptional} for $\std_v$ and
\eqref{eq:semiorthogonal} for $\boldsymbol\std^\sigma$.

\item\label{step:fullness}
\textbf{Fullness in $\D_{\fd}(A)$.}
Let $J\subseteq A$ be the ideal generated by all arrows.
Acyclicity implies $d_A(J)\subseteq J$.
For each vertex $v$, define the simple dg module
$S_v=P_v/e_vJ$, concentrated in degree $0$.
The outgoing arrows form a graded free basis of $e_vJ$ over the
corresponding projectives. The generator-degree argument of
\eqref{eq:generator-differential}--\eqref{eq:kernel-filtration}
therefore gives a finite projective filtration of $e_vJ$.
The exact sequence
$0\to e_vJ\to P_v\to S_v\to0$ makes every $S_v$ perfect.
The modules $S_1,\ldots,S_n$ are the simple $H^0(A)$-modules.
Finite-length filtrations and cohomological truncations give
\[
 \D_{\fd}(A)=\thick_{\D_{\fd}(A)}(S_1,\ldots,S_n)\subseteq\per A.
\]
Properness and \eqref{eq:perfect-generation} prove \eqref{eq:path-full}.

\item\label{step:pairing}
\textbf{The specified left-dual objects.}
For vertices $u,v$, restriction along $K_u\hookrightarrow P_u$ defines
a cochain map
\begin{equation}\label{eq:restriction}
 \rho_{uv}:\dgHom_A(P_u,\costd_v)
       \longrightarrow\dgHom_A(K_u,\costd_v).
\end{equation}
If $u<_\sigma v$, a basis of $e_uC_v$ consists of paths
$q:u\rightsquigarrow v$ all of whose vertices are
$\leq_\sigma v$. Each such $q$ has a unique factorisation
$q=ps$ with $p\in\mathcal F_u$.
Write $q^\vee$ and $s^\vee$ for the corresponding basis functionals.
The homogeneous $A$-linear map
$f_q\in\dgHom_A(P_u,\costd_v)^{-|q|}$ determined by
$f_q(e_u)=q^\vee$ satisfies
\begin{equation}\label{eq:restriction-basis}
 \rho_{uv}(f_q)(p')=
 \begin{cases}
 s^\vee,&p'=p,\\
 0,&p'\neq p
 \end{cases}
 \qquad(p'\in\mathcal F_u).
\end{equation}
Conversely, every generator-assignment basis vector of
$\dgHom_A(K_u,\costd_v)$ equals $\rho_{uv}(f_q)$ for a unique
path $q$ belonging to the stated basis of $e_uC_v$.
Hence $\rho_{uv}$ is an isomorphism of cochain complexes.

If $v<_\sigma u$, both complexes in \eqref{eq:restriction} vanish.
If $u=v$, the source of \eqref{eq:restriction} is $\Bbbk$ in degree $0$
and its target is zero. Applying $\RHom_A(-,\costd_v)$ to
\eqref{eq:path-short-exact}, using the homotopical projectivity of
$P_u$ and $K_u$, proves \eqref{eq:pairing}.

To verify exceptionality of
$(\costd_{\sigma(n)},\ldots,\costd_{\sigma(1)})$,
apply Steps~\ref{step:kernels}--\ref{step:fullness} to the graded
opposite algebra $A^{\op}$.
The left dg modules $C_v=B_ve_v$ are the resulting standard objects.
Thus $(C_{\sigma(1)},\ldots,C_{\sigma(n)})$ is exceptional in
$\D_{\fd}(A^{\op})$. Finite-dimensional vector-space duality gives
\[
 \Hom_{\D_{\fd}(A)}(\costd_u,\costd_v[t])
 \cong
 \Hom_{\D_{\fd}(A^{\op})}(C_v,C_u[t]).
\]
The sequence $(\costd_{\sigma(n)},\ldots,\costd_{\sigma(1)})$
is therefore exceptional.
Equation~\eqref{eq:path-full} places every $\costd_v$ in
$\thick_{\D_{\fd}(A)}(\std_1,\ldots,\std_n)$.

\item\label{step:amplitudes}
\textbf{Cohomological bounds and deflations.}
The modules $P_v,K_v,\std_v$, and $C_v$ have graded components only in
$[1-M,0]$, by \eqref{eq:path-width} and \eqref{eq:path-objects}.
Therefore
\[
 P_v,K_v,\std_v\in\H_A^M,\qquad
 \costd_v=\dual C_v\in\D_{\fd}(A)^{[0,M-1]}.
\]
Sequence \eqref{eq:path-short-exact} induces a conflation in $\H_A^M$,
so $p_v$ is a deflation. Conditions~\ref{ax:exceptional}--\ref{ax:full}
follow from Steps~\ref{step:exceptional}--\ref{step:amplitudes}.
For $m\geq M$, both cohomological windows only enlarge.
\end{enumerate}
\end{proof}

\begin{corollary}[Zero differential]\label{cor:zero-differential}
If $d_A=0$, then $M$ in \eqref{eq:path-width} is the smallest integer
for which $A$ satisfies \eqref{eq:standing}. Every ordering $\sigma$
therefore gives an $M$-reflexive qh dga.
\end{corollary}
\begin{proof}
For $d_A=0$, a path of minimal degree represents a non-zero class in
$H^{1-M}(A)$. Apply Theorem~\ref{thm:paths}.
\end{proof}

\begin{proposition}[Differentials on tree quivers]\label{prop:trees}
If the underlying unoriented graph of $Q$ is a tree, then the
differential $d_A$ in \eqref{eq:path-dga} is zero.
\end{proposition}
\begin{proof}
For an arrow $a:u\to v$, closure of $e_u,e_v$ gives
$d_A(a)\in e_uA^{|a|+1}e_v$. The only path from $u$ to $v$ is $a$,
so $e_uAe_v=\Bbbk a$ is concentrated in degree $|a|$.
Thus $d_A(a)=0$ for every arrow $a$; the derivation $d_A$ vanishes.
\end{proof}

\section{Examples}\label{sec:examples}

In Examples~\ref{ex:a4}--\ref{ex:differential}, the objects
$\std_v,K_v,C_v$, and $\costd_v$ are defined by
\eqref{eq:path-objects}. The symbol $S_v$ denotes the simple module
at vertex $v$, as in Step~\ref{step:fullness} of the proof of
Theorem~\ref{thm:paths}. All quotient modules inherit the indicated
path actions.

\begin{example}[A graded quiver of type $A_4$]\label{ex:a4}
Let
\[
 Q:\ 1\xrightarrow{a}2\xrightarrow{b}3\xrightarrow{c}4,
 \qquad |a|=|b|=|c|=-1,\qquad d_A=0.
\]
The path $abc$ has degree $-3$, so $M=4$.
Choose $2<_\sigma4<_\sigma1<_\sigma3$, and put
$U=P_1/(abA)$, with basis $e_1,a$. Formula~\eqref{eq:path-objects} gives
\begin{equation}\label{eq:a4-objects}
\begin{array}{c|c|c|c}
 v&\std_v&K_v&\costd_v\\ \hline
 1&U&abA\cong P_3[2]&S_1\\
 2&S_2&bA\cong P_3[1]&S_2\\
 3&P_3&0&\dual(Ae_3)\\
 4&S_4&0&S_4
\end{array}
\end{equation}
The module $Ae_3$ has basis $e_3,b,ab$ in degrees $0,-1,-2$.
Hence
\begin{equation}\label{eq:a4-cohomology}
 \begin{gathered}
 H^0(U)=S_1,\quad H^{-1}(U)=S_2,\qquad
 H^0(P_3)=S_3,\quad H^{-1}(P_3)=S_4,\\
 H^0(\costd_3)=S_3,\quad H^1(\costd_3)=S_2,\quad
 H^2(\costd_3)=S_1.
 \end{gathered}
\end{equation}
Theorem~\ref{thm:paths} gives the full exceptional sequence
$(S_2,S_4,U,P_3)$ with left-dual sequence
$(\costd_3,S_1,S_4,S_2)$.
The presentations of $U$ and $S_2$ are the conflations
\[
 P_3[2]\xrightarrow{x\mapsto abx}P_1\longrightarrow U\dashrightarrow,
 \qquad
 P_3[1]\xrightarrow{x\mapsto bx}P_2\longrightarrow S_2\dashrightarrow
 \quad\text{in }\H_A^4.
\]
\end{example}

\begin{example}[A graded quiver of type $D_4$]\label{ex:d4}
Let
\[
 Q:\ 1\xrightarrow{a}2\xrightarrow{b}3,\quad 2\xrightarrow{c}4,
 \qquad |a|=|c|=-1,\quad |b|=0,\quad d_A=0.
\]
The path $ac$ has degree $-2$, so $M=3$.
Choose $1<_\sigma3<_\sigma2<_\sigma4$, and set
$V=P_2/(cA)$, with basis $e_2,b$ in degree $0$. Then
\begin{equation}\label{eq:d4-objects}
\begin{array}{c|c|c|c}
 v&\std_v&K_v&\costd_v\\ \hline
 1&S_1&aA\cong P_2[1]&S_1\\
 2&V&cA\cong P_4[1]&\dual(Ae_2)\\
 3&S_3&0&S_3\\
 4&S_4&0&\dual(Ae_4)
\end{array}
\end{equation}
The left modules $Ae_2$ and $Ae_4$ have respective bases
$e_2,a$ and $e_4,c,ac$. Consequently,
\begin{equation}\label{eq:d4-cohomology}
 \begin{gathered}
 H^0(\costd_2)=S_2,\quad H^1(\costd_2)=S_1,\\
 H^0(\costd_4)=S_4,\quad H^1(\costd_4)=S_2,\quad
 H^2(\costd_4)=S_1.
 \end{gathered}
\end{equation}
The standard sequence $(S_1,S_3,V,S_4)$ and its left-dual sequence
$(\costd_4,\costd_2,S_3,S_1)$ satisfy
Definition~\ref{def:reflexive} with $m=3$.
For example,
$H^{-1}(K_1)=V$ and $H^{-2}(K_1)=S_4$, so the conflation
$K_1\to P_1\to S_1\dashrightarrow$ belongs to $\H_A^3$.
\end{example}

\begin{example}[A non-zero differential and the choice of width]
\label{ex:differential}
Take
\begin{equation}\label{eq:differential-example}
 \begin{gathered}
 Q:\ 1\xrightarrow{a}2\xrightarrow{b}3\xrightarrow{h}4,
 \qquad c:1\to3,\\
 |a|=|b|=0,\quad |c|=|h|=-1,\qquad
 d_A(c)=ab,\quad d_A(a)=d_A(b)=d_A(h)=0.
 \end{gathered}
\end{equation}
The only non-zero path differentials are $d_A(c)=ab$ and
$d_A(ch)=abh$. Thus
\begin{equation}\label{eq:differential-cohomology}
 H^0(A)=\Bbbk\{e_1,e_2,e_3,e_4,a,b\},\qquad
 H^{-1}(A)=\Bbbk\{h,bh\},
 \qquad H^q(A)=0\ (q\notin[-1,0]).
\end{equation}
The smallest parameter allowed by \eqref{eq:standing} is $m=2$,
whereas the path $ch$ gives $M=3$ in \eqref{eq:path-width}.

Choose $3<_\sigma4<_\sigma1<_\sigma2$, and put
$X=P_1/(aA)$. The quotient $X$ has basis $e_1,c,ch$ and zero
differential. Formula~\eqref{eq:path-objects} gives
\begin{equation}\label{eq:differential-standards}
 \begin{gathered}
 (\std_{\sigma(1)},\std_{\sigma(2)},\std_{\sigma(3)},\std_{\sigma(4)})
       =(S_3,S_4,X,P_2),\\
 (K_1,K_2,K_3,K_4)=(aA,0,hA,0),\\
 H^0(X)=S_1,\qquad H^{-1}(X)=S_3,\qquad H^{-2}(X)=S_4.
 \end{gathered}
\end{equation}
The objects in \eqref{eq:differential-standards}, with the objects
$\costd_v$ of \eqref{eq:path-objects}, define a $3$-reflexive qh dga
by Theorem~\ref{thm:paths}. The same standard $X$ does not belong to
$\H_A^2$, because the class of $ch$ generates $H^{-2}(X)$.

A different order isolates the condition on the left-dual objects.
Choose $1<_\tau3<_\tau4<_\tau2$, and denote the objects
of \eqref{eq:path-objects} for $\tau$ by
$\std_v^\tau,K_v^\tau,C_v^\tau$, and $\costd_v^\tau$.
Then
\begin{equation}\label{eq:differential-costandards}
 \begin{gathered}
 (\std_1^\tau,\std_2^\tau,\std_3^\tau,\std_4^\tau)
       =(S_1,P_2,S_3,S_4),\\
 C_4^\tau=(A/Ae_2A)e_4=\Bbbk\{e_4,h,ch\},
 \qquad d_{C_4^\tau}=0,\\
 H^0(\costd_4^\tau)=S_4,\quad H^1(\costd_4^\tau)=S_3,\quad
 H^2(\costd_4^\tau)=S_1.
 \end{gathered}
\end{equation}
Every $\std_v^\tau$ and $K_v^\tau$ belongs to $\H_A^2$.
Indeed, $K_1^\tau=\Bbbk\{a,ab,abh,c,ch\}$ has cohomology $S_2$
in degree $0$, $K_3^\tau=hA\cong P_4[1]$, and
$K_2^\tau=K_4^\tau=0$.
However, $\costd_4^\tau$ in \eqref{eq:differential-costandards} violates
\eqref{eq:costandard-window} for $m=2$.
Theorem~\ref{thm:paths} applies to $\tau$ with $m=3$.

Equations~\eqref{eq:differential-standards} and
\eqref{eq:differential-costandards} concern the specific construction
\eqref{eq:path-objects}; they do not exclude other standard families
for $\sigma$ or $\tau$ at width $2$.
\end{example}

\section{Questions}\label{sec:questions}

Theorem~\ref{thm:paths} settles existence for graded path dg algebras
at every width $m\geq M$. For a general algebra satisfying
\eqref{eq:standing}, several structural questions remain.

\subsection*{Foundational questions}
\begin{enumerate}[label=\textbf{F\arabic*.}]
 \item \textbf{Generation and perfection.}
 If conditions~\ref{ax:exceptional}--\ref{ax:costandard} hold, when does
 \eqref{eq:full} follow? For an $m$-reflexive qh dga, when do the
 standard objects $\std_v$ belong to $\per A$?
 \item \textbf{A prescribed cohomological width.}
 Characterise the orders admitting Definition~\ref{def:reflexive}
 for a fixed $m$ satisfying \eqref{eq:standing}.
 In Example~\ref{ex:differential}, determine whether standard families
 other than \eqref{eq:path-objects} exist at width $2$ for the orders
 $\sigma$ and $\tau$.
 \item \textbf{Filtrations by shifted standards.}
 Determine hypotheses under which the presentations
 \eqref{eq:presentation} can be chosen with
 \begin{equation}\label{eq:shifted-question}
 K_v\in\Filt_{\H_A^m}
 \{\std_w[s]:v<_\sigma w,\ 0\leq s<m,\ \std_w[s]\in\H_A^m\}.
 \end{equation}
 Establish the relation between \eqref{eq:shifted-question} and the
 cohomological requirement \eqref{eq:costandard-window}.
\end{enumerate}

\subsection*{Further directions}
\begin{enumerate}[label=\textbf{P\arabic*.}]
 \item \textbf{Heredity and recollement.}
 Relate Definition~\ref{def:reflexive} to filtrations by dg ideals and
 recollements whose successive factors are equivalent to
 $\D^b(\mod\Bbbk)$.
 \item \textbf{Standardisation and Ringel duality.}
 Study the extension-closed subcategories
 $\Filt_{\H_A^m}\{\std_v\}$ and
 $\Filt_{\H_A^m}\{\costd_v[m-1]\}$.
 Determine conditions producing a tilting or silting object $T$ and an
 $m'$-reflexive qh structure on $\REnd_A(T)$; determine the resulting
 integer $m'\geq1$.
 \item \textbf{Reciprocity and dimension bounds.}
 Seek a graded form of BGG reciprocity and homological bounds in terms
 of $m$ and $n$. In the Grothendieck group $K_0(\D_{\fd}(A))$, write
 $[Y]$ for the class of $Y\in\D_{\fd}(A)$. Filtrations in
 \eqref{eq:shifted-question} must account for
 $[Y[s]]=(-1)^s[Y]$ for $s\in\ZZ$.
\end{enumerate}

\begin{thebibliography}{MP25}
\bibitem[BB26]{BB26}
A. Bodzenta and A. Bondal,
\emph{Highest weight categories via pairs of dual exceptional sequences},
2026, \href{https://arxiv.org/abs/2601.22004}{arXiv:2601.22004}.

\bibitem[DR92]{DR92}
V. Dlab and C. M. Ringel,
\emph{The module theoretical approach to quasi-hereditary algebras},
in \emph{Representations of algebras and related topics},
London Mathematical Society Lecture Note Series 168,
Cambridge University Press, 1992, pp.~200--224.
\href{https://www.math.uni-bielefeld.de/~ringel/opus/good-d-r.pdf}{Author's manuscript}.

\bibitem[MP25]{MP25}
N. Mochizuki and M. Plogmann,
\emph{On the Auslander--Reiten theory for extended hearts of proper
connective dg algebras},
2025, \href{https://arxiv.org/abs/2505.16560}{arXiv:2505.16560}.

\end{thebibliography}
\end{document}
